% sec_experiments_journal.tex
% Journal version: appendix content merged in-text, two-column layout fixes.

\section{Experiments}
\label{sec:experiments}

\subsection{Experimental Setup}

\paragraph{Datasets.}
\textbf{DAD}~\citep{ref4}: 1,750 dashcam videos (620 accident,
1,130 non-accident) from six Taiwanese cities, split 1,280/450
train/test at 20\,fps (100 frames).
\textbf{A3D}~\citep{liao2024crash}: 1,199 videos (960/239) at 10\,fps,
100 frames, diverse East Asian urban accident types.

\paragraph{Metrics.}
\textbf{AP}: area under the precision-recall curve.
\textbf{mTTA}: mean lead time at first issued alert.
\textbf{TTA@R80} / \textbf{P@R80}: TTA and precision at 80\% recall.

\paragraph{Input representation and concept vocabulary.}
Each frame provides $N{=}19$ object-centric VGG16~\citep{ref31} ROI
features of dimension $D{=}4096$.
The concept bottleneck uses $K{=}837$ traffic-relevant concepts spanning
agent behaviour, spatial risk, road condition, and environmental context,
derived from a large candidate set filtered for semantic plausibility.

\paragraph{Optimisation.}
AdamW, lr $3 \times 10^{-4}$ with cosine decay to $10^{-5}$, batch 16,
$h=256$, $z=128$. Training uses staged scheduling: classification and
concept supervision dominate early epochs; Actor--Critic grows later.
For DAD, a \emph{curriculum-stabilised} variant trains the backbone
for 20 epochs without CBM loss, then ramps concept losses over 30 epochs.

\subsection{Interpretability Taxonomy}

\begin{table*}[t]
\caption{Interpretability comparison. \textbf{WHY}: intrinsic concept
explanation. \textbf{WHEN}: intrinsic alert-timing explanation.
\method{} provides both within a single model.}
\label{tab:interp_comparison}
\centering\small
\begin{tabular}{lcccc}
\toprule
Method & WHY & WHEN & Intrinsic? & Params\\
\midrule
DSA-RNN~\citep{ref4}         & \texttimes & \texttimes & -- & 12M\\
DRIVE~\citep{ref2}           & Grad-CAM  & \texttimes & post-hoc & 31M\\
DSTA~\citep{ref17}           & \texttimes & \texttimes & -- & 180M\\
DAA-GNN~\citep{ref32}        & \texttimes & \texttimes & -- & 183M\\
W3AL~\citep{liao2024and}     & \texttimes & verbal    & post-hoc & 7B\\
CRASH~\citep{liao2024crash}  & \texttimes & \texttimes & -- & 26M\\
CCAF-Net~\citep{liu2025ccaf} & \texttimes & \texttimes & -- & 191M\\
LATTE~\citep{ZHANG2025103173}& \texttimes & \texttimes & -- & 45M\\
\midrule
\textbf{\method{} (Ours)} & \textbf{837 concepts} &
  \textbf{AC policy} & \textbf{\checkmark} & \textbf{30M}\\
\bottomrule
\end{tabular}
\end{table*}

Table~\ref{tab:interp_comparison} shows a clear gap: no prior method provides
both WHY and WHEN interpretability intrinsically.
DRIVE's Grad-CAM is post-hoc and carries no concept structure;
W3AL requires a 7B LLM for verbal localisation.
\method{} closes both gaps within only 30M parameters.

\subsection{Why the Dual-Layer Decomposition Is Well Motivated}

A conventional recurrent anticipation model predicts risk from a single latent
state $\mathbf{h}_t$, offering no separation between evidence extraction and
decision timing. \method{} decomposes these roles: the WHY layer extracts
$\mathbf{c}_t$, and the WHEN layer learns a timing policy over
$\mathbf{s}_t=[\mathbf{h}_t\|\mathbf{c}_t]$.

This changes the information geometry of the problem. If a hazard concept rises
before impact, the policy can cross its alert boundary before the latent
classifier reaches peak confidence---the concept channel acts as an
\emph{early-warning basis}. Furthermore, a perturbation
$\Delta\mathbf{c}_t$ to a concept in the pre-crash window induces
$\mathbf{s}'_t=[\mathbf{h}_t\|\mathbf{c}_t+\Delta\mathbf{c}_t]$, so
semantics have \emph{direct causal leverage} on timing decisions rather than
being computed post hoc. Section~\ref{sec:intervention} provides empirical
evidence for this claim.
